\subsection{Every rational toric pair in the inherited cubic class-two quotient}

\begin{statement}{E.29}{Proved}{The exact tame-residue criterion and the three classes in each nonempty fibre}
Let $F$ be the free pro-$3$ geometric fundamental group of
$\mathbb P^1-\{0,1,\infty\}$, based at the rational tangent
$t=0$, $dt=1$, and let $N_3=\ker(F\to(\mathbb Z/3)^2)$.
The image of $N_3$ in $F/\overline{\gamma_3F}$ is the pro-$3$
group $H$ with coordinates $(u,v,w)\in\mathbb Z_3^3$, law
\[
 (u,v,w)(u',v',w')
 =\left(u+u',v+v',w+w'+\frac92(uv'-vu')\right),              \tag{E.29a}
\]
and arithmetic action
\[
 g(u,v,w)=(\chi_3(g)u,\chi_3(g)v,\chi_3(g)^2w).
\]
It is the inherited class-two quotient of the plus cubic;
its arithmetic extension, with the lifted tangent, is
$H\rtimes G_{\mathbb Q}$.

For $a,b\in\mathbb Q^*$ let $\mathscr F(a,b)$ be the fibre over
$(\kappa_a,\kappa_b)$ of the continuous section-class map
\[
 H^1_{\mathrm{cts}}(G_{\mathbb Q},H)
 \longrightarrow H^1_{\mathrm{cts}}
       (G_{\mathbb Q},\mathbb Z_3(1)^2).
\]
For a rational prime $q$, put
\[
 \theta_q(a,b)=
 (-1)^{v_q(a)v_q(b)}
 \overline{a^{v_q(b)}/b^{v_q(a)}}\ \in\mathbb F_q^*.
                                                               \tag{E.29b}
\]
The expression under the bar is a $q$-adic unit. Let
$o_3(\theta)$ denote the largest power of three dividing its
multiplicative order. Then the complete answer on rational
toric pairs is
\[
 |\mathscr F(a,b)|=
 \begin{cases}
 3,&o_3(\theta_q(a,b))\le9\text{ for every prime }q,\\
 0,&o_3(\theta_q(a,b))>9\text{ for at least one prime }q.
 \end{cases}                                                   \tag{E.29c}
\]
Only primes dividing a numerator or denominator of $a$ or $b$
need to be tested. For each nonempty fibre, choose one continuous
cocycle $z(g)\in H$. Its three distinct classes are represented by
\[
 z_j(g)=\left(0,0,j\frac{\chi_3(g)^2-1}{3}\right)z(g),
 \qquad j=0,1,2.                                               \tag{E.29d}
\]
The choice of $z$ is the stated ambiguity in this parametrization.
The obstruction itself is exactly
$9\kappa_a\cup\kappa_b\in
H^2_{\mathrm{cts}}(G_{\mathbb Q},\mathbb Z_3(2))$.

For example, $(6,109)$ and $(2,163)$ have empty fibres,
whereas $(6^3,109)$ and $(2^9,163)$ each have three classes.
The source of this classification is the displayed inherited
quotient. The intrinsic rank-ten class-two quotient maps onto
it by
\[
 N_3/\overline{\gamma_3N_3}\twoheadrightarrow H,
 \qquad
 \ker=(N_3\cap\overline{\gamma_3F})/
                     \overline{\gamma_3N_3}.                  \tag{E.29e}
\]
Thus its section classes map to the fibres computed here.

\emph{Manuscript: \S51.16, folios 733v--736v, scans 172--178,
(371)--(379); the Kummer coordinates of the minimum quotient
are retained, and their next inherited integral obstruction is
computed. This statement treats every rational toric pair.}
\end{statement}
\begin{proof}
We first fix the comparison and its factors of three. Put
$X=\log x$, $Y=\log y$, and $C=[X,Y]$ in the integral
class-two pro-$3$ group; the denominator $2$ in its
Baker--Campbell--Hausdorff law is a unit in $\mathbb Z_3$.
The image of $N_3$ has the lattice
\[
 A=3X,\qquad B=3Y,\qquad C=[X,Y],\qquad[A,B]=9C.
\]
Indeed its abelian coordinates are divisible by three,
and $x^3,y^3,[x,y]\in N_3$ generate all these coordinates.
This proves (E.29a) and the exact surjection (E.29e).
For the normalized arithmetic action,
$x\mapsto x^{\chi_3(g)}$ and
$y\mapsto f_g^{-1}y^{\chi_3(g)}f_g$, where $f_g\in F'$.
The element $f_g$ is central after quotienting by
$\overline{\gamma_3F}$, giving the displayed diagonal action.

On the plus cubic put $u=x_0/x_\infty$, $v=x_1/x_\infty$.
The covering is $t=-u^3$, $1-t=-v^3$.
At its rational cusp above $t=0$, the parameter $s=-u$ gives
$t=s^3$ and $v=-(1-s^3)^{1/3}$ with constant coefficient $-1$.
The tangential Galois section at $s=0$, $ds=1$ maps to the
chosen base section: the corresponding embedding of Puiseux
fields sends $t^{1/m}$ to $s^{3/m}$ and commutes with
coefficientwise $G_{\mathbb Q}$-action. This specifies the
comparison at the ramified cusp.
Both leading toric units are $-1$, whose pro-$3$ Kummer
classes are zero, since $-1$ has compatible rational
$3^n$th roots. Thus a toric cocycle $(\kappa_a,\kappa_b)$
has inherited coordinates $\kappa_a A+\kappa_b B$ exactly.
The comparison to base toric coordinates is multiplication
by three, coming from $t=-u^3$, $1-t=-v^3$.

Write $V=\mathbb Z_3(1)^2$, $Z=\mathbb Z_3(2)$, so that
$1\to Z\to H\to V\to1$ is a central extension of pro-$3$
groups with $G_{\mathbb Q}$-action. For a continuous
$V$-cocycle $v_g=(\alpha_g,\beta_g)$, use its continuous
set-theoretic lift $(\alpha_g,\beta_g,0)$. Its obstruction
two-cocycle is
\[
 \omega(g,h)=\frac92\left(
 \alpha_g\chi_3(g)\beta_h-
 \beta_g\chi_3(g)\alpha_h\right).                           \tag{E.29f}
\]
The symmetric sum of the two cup cocycles is the negative
coboundary of $g\mapsto\alpha_g\beta_g$. Hence the class
of (E.29f) is $9[\alpha\cup\beta]$.
Vanishing of this class means precisely that a continuous
central one-cochain cancels (E.29f), which is the cocycle
condition for a lift to $H$.

We now prove the criterion in continuous cohomology itself.
The pairing
\[
 (a,b)\longmapsto\kappa_a\cup\kappa_b
 \quad\text{in }H^2_{\mathrm{cts}}(G_{\mathbb Q},Z)
\]
is bilinear and satisfies the integral Steinberg relation.
For the latter, the rational point $t=a$ of
$\mathbb P^1-\{0,1,\infty\}$, for $a\ne0,1$, supplies
an actual arithmetic section. Its class-two pro-$3$
quotient has toric cocycle $(\kappa_a,\kappa_{1-a})$ and
commutator coefficient one. The calculation (E.29f) with
$9$ replaced by $1$ therefore gives
$\kappa_a\cup\kappa_{1-a}=0$ in continuous integral
cohomology. Bilinearity now defines a homomorphism
\[
 \beta:K_2^M(\mathbb Q)\longrightarrow
 H^2_{\mathrm{cts}}(G_{\mathbb Q},Z),\qquad
 \{a,b\}\longmapsto\kappa_a\cup\kappa_b.                    \tag{E.29g}
\]
This constructs (E.29g) before taking any inverse limit of
degree-two cohomology groups.

We use Matsumoto's presentation and the localization
calculation in the opened source
\cite[III.6.1, III.6.5--III.6.5.1, III.5.2.2]{E_Kbook}.
For $\mathbb Q$ these give the exact sequence
\[
 0\longrightarrow\mathbb Z/2\longrightarrow K_2^M(\mathbb Q)
 \xrightarrow{(\partial_q)_q}
 \bigoplus_q\mathbb F_q^*\longrightarrow0.                  \tag{E.29h}
\]
The kernel is generated by $\{-1,-1\}$. In the convention
of that source, $\partial_q\{a,b\}=\theta_q(a,b)^{-1}$;
coordinatewise inversion is the explicit automorphism of
the right-hand side relating the two conventions.
Every element of the middle group is torsion: its finitely
supported image has finite order, and its remaining kernel
is killed by two. In particular (E.29h) identifies its
$3$-primary subgroup with
$\bigoplus_q(\mathbb F_q^*)\{3\}$.
The target of (E.29g) is a $\mathbb Z_3$-module, so
$\beta$ kills all primary components at primes different
from three.

We give the local detection argument to prove injectivity
on the remaining $3$-primary subgroup. Let $q\ne3$.
For $T_n=\mu_{3^n}^{\otimes2}$, wild inertia has no positive
cohomology with these coefficients, and the $3$-part of
tame inertia is $\mathbb Z_3(1)$. The quotient Galois group
of the finite residue field and this tame $3$-part have
$3$-cohomological dimension one. Their two-term procyclic
resolutions, followed by Hochschild--Serre, give
\[
 H^2(\mathbb Q_q,T_n)
  \simeq H^1(\mathbb F_q,\mu_{3^n})
  \simeq (\mathbb Z/3^n)/(q-1).                              \tag{E.29i}
\]
In particular the degree-one finite-coefficient local
groups are finite, so their first derived inverse limit
vanishes. Passing to the inverse limit of (E.29i) therefore
gives an actual continuous-cohomology isomorphism
\[
 \operatorname{res}^{\mathrm{tame}}_q:
 H^2_{\mathrm{cts}}(\mathbb Q_q,Z)
   \xrightarrow{\sim}
 H^1_{\mathrm{cts}}(\mathbb F_q,\mathbb Z_3(1))
   \simeq\mathbb Z_3/(q-1).                                 \tag{E.29j}
\]

Here is the cup-product normalization of this map.
Choose arithmetic Frobenius $\phi$ and a tame generator
$\tau$ with $\phi\tau\phi^{-1}=\tau^q$; these are the
operators on compatible roots of a uniformizer. The
conjugation formula is also given in the opened source
\cite[(8.2.2)]{E_BrinonConrad}.
For a unit $u$, its Kummer class is unramified. Write
$a=\kappa_u(\phi)$ and choose $\phi$ fixing the chosen
uniformizer roots, so $\kappa_q(\phi)=0$ and
$\kappa_q(\tau)=1$. The extension two-cocycle
$\kappa_u\cup\kappa_q$ has the relation defect $qa$
between $\phi\tau\phi^{-1}$ and $\tau^q$.
Changing the central lift of $\tau$ by $b$ changes that
defect by $q(q-1)b$. Thus (E.29j), which divides the
defect by $q$, sends this cup product to the unit's
residue-field Kummer class. This specifies the map and its
sign, and verifies the edge-map calculation directly.
Two unit Kummer classes have zero cup product because
they factor through the procyclic residue Galois group.
Also $\kappa_q\cup\kappa_q=0$, since graded commutativity
kills it by two and two is a unit in $\mathbb Z_3$.
Writing $a=q^\alpha u$, $b=q^\beta v$ now yields
\[
 \operatorname{res}^{\mathrm{tame}}_q
    (\kappa_a\cup\kappa_b)
    =\kappa_{\overline{u}^{\beta}/\overline{v}^{\alpha}}
    =\kappa_{\theta_q(a,b)}.                                 \tag{E.29k}
\]
The sign in (E.29b) has zero pro-$3$ Kummer class.
Finite-field Kummer theory identifies the image of an
element of $\mathbb F_q^*$ in (E.29j) with its
$3$-primary component, preserving its order.

Every nonzero $3$-primary element of $K_2^M(\mathbb Q)$
has a nonzero tame component at some $q\ne3$, by
(E.29h), since $\mathbb F_3^*$ has order two.
Equation (E.29k) detects that component under $\beta$.
Consequently $\beta$ is injective on this subgroup.
It follows that
\[
 9\kappa_a\cup\kappa_b=0
 \quad\Longleftrightarrow\quad
 9\{a,b\}_{\{3\}}=0
 \quad\Longleftrightarrow\quad
 o_3(\theta_q(a,b))\le9\quad\text{for every }q.               \tag{E.29l}
\]
This proves both necessity and sufficiency in continuous
cohomology, including sufficiency for the existence of a
continuous central cochain. No global inverse-limit
injectivity assertion was used.

It remains to compute each nonempty fibre. We first compute
\[
 H^1_{\mathrm{cts}}(G_{\mathbb Q},\mathbb Z_3(2))
      \simeq\mathbb Z/3.                                     \tag{E.29m}
\]
A continuous cocycle with values in $\mathbb Q_3(2)$
restricts on inertia at $q\ne3$ to a homomorphism.
Wild inertia contributes zero. If $a$ is its value on a
tame $3$-generator, Frobenius equivariance gives
$q^2a=qa$, so $(q-1)a=0$, hence $a=0$.
It therefore vanishes on the closed normal subgroup
generated by these inertia groups and factors through
$G_{\mathbb Q,\{3\}}$.
Theorem 7.2 and Theorem B.1 of the opened
Hain--Matsumoto paper \cite{E_HainMatsumoto} give
\[
 \dim_{\mathbb Q_3}H^1_{\mathrm{cts}}
     (G_{\mathbb Q,\{3\}},\mathbb Q_3(2))=r_2(\mathbb Q)=0.
\]
This proves the vanishing over the full $G_{\mathbb Q}$.
The exact coefficient sequence
$0\to\mathbb Z_3(2)\to\mathbb Q_3(2)
\to\mathbb Q_3/\mathbb Z_3(2)\to0$ thus identifies (E.29m)
with $H^0(G_{\mathbb Q},\mathbb Q_3/\mathbb Z_3(2))$.
The image of $\chi_3$ is $\mathbb Z_3^*$; the ideal generated
by $u^2-1$ for all units $u$ is $3\mathbb Z_3$.
The invariants are consequently
$(\tfrac13\mathbb Z_3)/\mathbb Z_3$, and the boundary of
$j/3$ is exactly the central cocycle in (E.29d).

To see that this calculation gives the whole fibre, choose
one lifted cocycle. The ratio of two lifts with the same
toric cocycle is a central one-cocycle, and every such
cocycle gives another lift. A kernel conjugation which
preserves the exact toric cocycle has projection in
$H^0(G_{\mathbb Q},V)=0$: complex conjugation acts by
$-1$ on the torsion-free group $V$. Such a conjugation
is therefore central and changes the ratio only by a
central coboundary. When the toric cocycles initially
differ by a coboundary, the surjection $H\to V$ supplies
the conjugation which makes them equal. Thus the fibre
is a torsor under the explicitly computed group (E.29m),
with exactly the representatives (E.29d).

The program \texttt{code/e\_rational\_pairs.py} computes
(E.29b) by exact rational arithmetic. Among all $40\,000$
ordered positive integer pairs up to $200$, it finds
$39\,378$ nonempty and $622$ empty fibres. The histogram
of the largest $3$-primary residue order is
\[
\begin{array}{c|rrrrr}
 \text{order}&1&3&9&27&81\\\hline
 \text{number of pairs}&19036&15198&5144&364&258.
\end{array}
\]
It also checks all $101\,124$ ordered pairs of the $318$
nonzero signed reduced rationals of numerator and
denominator height at most $16$, and all $317$ applicable
Steinberg pairs in that set. It checks (E.29d) as a
cocycle modulo $3^n$ for $1\le n\le4$, and the three
twist-two invariants and absence of twist-one invariants
at those finite levels. All output and every obstructed
integer pair are retained in the accompanying JSON file.
\end{proof}
